\section{Discussion}
\label{sec:discussion}

\paragraph{What limits the model.}
At matched width the causal PT's readout has a transformer's bottleneck and its state has the
same width, so what remains is the dynamics. The
richer exact readout trains worse (\Cref{sec:exp3}), and in a one-seed sweep more iterations did
not help (\Cref{sec:exp-depth}). This suggests that optimisation, not expressivity, is the
limit at this scale. The exact readout's penalty falls as the label variable is split
(\Cref{sec:exp-factored}): with one label variable it may have had little extra to express.

\paragraph{Tuning and the feed-forward block.}
Much of the width sensitivity has one mechanism, the head temperature; the ceiling that remains
after correcting it is unexplained (see Limitations). The failures at a large $L_2$
coefficient, at $\clip=7$ and under the original temperature share one cause: the ratio of the
arc message to the word unary leaves the range in which the label update is stable
(\Cref{app:res-knobs}); at $\clip=0$ the model cannot see distance. This ratio is worth
monitoring while tuning, because perplexity stays flat until it jumps. The head temperature,
learning rate, rank and label-set size matter most, and fewer iterations are better (an ordering
from different budgets, not a ranking).
The feed-forward block is the costliest single
difference on the switch ladder that still trains, and the global variables proposed as its
in-graph substitute do not close the gap because their posterior stays near uniform, a symmetry
of the parameterisation that does not rule out other substitutes (\Cref{app:grid-globalhead}).
